\cvheader{Curriculum vitae}

\cvsection{Professional Experience}
\begin{cvchronology}
\cventry{2020--}{Professor at the School of Mathematics and Statistics of the University of Melbourne (Australia).}
\cventry{2016--2020}{ARC Future Fellow at the School of Mathematics and Statistics of the University of Melbourne (Australia).}
\cventry{2013--2016}{Directeur de recherches (senior scientist) CNRS at LPTHE (UMR 7589), Universit\'e Pierre et Marie Curie (France).}
\cventry{2012}{``Research Member'' of the MSRI program ``Random Spatial Processes'' (Berkeley).}
\cventry{2009--2010}{``Visiting Professor'' of the KNAW (Royal Netherlands Academy of Arts and Sciences), visiting the University of Amsterdam.}
\cventry{2008--2013}{Charg\'e de recherches CNRS at LPTHE (UMR 7589), Universit\'e Pierre et Marie Curie (France).}
\cventry{2005--2008}{Charg\'e de recherches CNRS at LPTMS (UMR 8626), Universit\'e Paris-Sud (France).}
\cventry{2004--2005}{Charg\'e de recherches CNRS at Laboratoire Poncelet (LIFR--MIIP, UMI 2615), Moscow (Russia).}
\cventry{2000--2004}{Charg\'e de recherches (junior scientist) CNRS at LPTMS (UMR 8626), Universit\'e Paris-Sud (France).}
\cventry{2000}{Visitor at C.N.~Yang Institute for Theoretical Physics, SUNY Stony Brook.}
\cventry{1999--2000}{Post-doctoral fellow at the New High Energy Theory Center, Rutgers University.}
\cventry{1998--1999}{National service.}
\cventry{1995--1998}{PhD in theoretical physics at the Laboratoire de Physique Th\'eorique de l'Ecole Normale Sup\'erieure de Paris, supervisor V.~Kazakov.}
\end{cvchronology}

\cvsection{Research Interests}
\begin{cvcompact}
\item Algebraic combinatorics, symmetric functions; Schubert calculus, enumerative geometry.
\item Exactly solvable models in statistical mechanics, quantum integrability, Yang--Baxter equation, and quantum groups.
\item Random matrices, and connections between mathematical physics and enumerative combinatorics.
\item Computational and experimental mathematics, including algorithm design, computer algebra and machine learning.
\end{cvcompact}

\cvsection{Prizes, Awards}
\begin{cvchronology}
\cventry{2016--2020}{Australian Research Council Future Fellowship ``From quantum integrable systems to algebraic geometry and combinatorics''.}
\cventry{2011--2016}{European Research Council grant ``Loop models, Integrability and Combinatorics''.}
\cventry{2010--2014}{Prime d'excellence scientifique (scientific excellency award) from CNRS.}
\cventry{2009}{Visiting Professorship award from the Royal Netherlands Academy of Arts and Sciences.}
%\cventry{2003}{President of the French Republic's award ``Creators/Innovators of the year 2003''.}
\end{cvchronology}

\cvsection{Teaching Experience}
\begin{cvchronology}
\cventry{2019--2026}{\textit{Discrete Mathematics}, \textit{Linear Algebra}, \textit{Group theory and linear algebra}, \textit{Enumerative combinatorics}, \textit{Exactly solvable models}, \textit{Experimental mathematics}\/ (48h/36h) courses at the University of Melbourne.}
\cventry{2022}{LAWRGe lectures at the University of Southern California: \textit{Schubert Calculus and Quantum Integrability} (14h).}
\cventry{2019}{Mini-course at the University of Virginia: \textit{Quantum integrability and symmetric polynomials} (10h).}
\cventry{2019}{Mini-course at the University of Melbourne: \textit{Schubert calculus} (12h).}
\cventry{2018}{Mini-course at the University of Melbourne: \textit{Geometry, Integrability and symmetric functions} (6h).}
\cventry{2015}{Lectures at the Department of Mathematics of the Higher School of Economics (Moscow): \textit{Geometry, Integrability and symmetric functions} (8h).}
\cventry{2013--2015}{\textit{Int\'egrabilit\'e quantique}, Masters course at Universit\'e Pierre et Marie Curie (30h).}
\cventry{2012}{Chern--Simons lectures at the University of California at Berkeley: \textit{Exactly solvable tiling models, Schur functions and the Littlewood--Richardson rule} (6h).}
\cventry{2008}{Course at Les Houches summer school: \textit{Integrability and combinatorics: selected topics}.}
\cventry{2007}{Lectures at Barcelona University: \textit{Quantum integrability and Razumov--Stroganov type conjectures}.}
\cventry{1997--1998}{Teaching Assistant at ENS Paris: mathematics course.}
\cventry{1995--1997}{Oral examiner at preparatory school Louis-Le-Grand (bachelor level).}
\end{cvchronology}

\cvsection{Studies and Academic Honours}
\begin{cvchronology}
\cventry{2008}{Habilitation \`a Diriger des Recherches ``Six-Vertex, Loop and Tiling models: Integrability and Combinatorics'' (Universit\'e Pierre et Marie Curie, Paris 6).}
\cventry{1998}{PhD ``Quelques applications de l'Ansatz de Bethe'', mention tr\`es honorable avec les f\'elicitations du jury (highest mark) (Universit\'e Pierre et Marie Curie, Paris 6).}
\cventry{1996}{Mathematics ``Agr\'egation'' (teaching diploma).}
\cventry{1994--1995}{D.E.A. (Master's) in theoretical physics, ENS.}
\cventry{1993--1994}{Licence et ma\^\i trise in mathematics, licence in physics (Bachelor's degrees).}
\cventry{1993}{Admission at Ecole Normale Sup\'erieure (ENS, Paris) on the Mathematics--Physics--Computer~Science entrance exam.}
\cventry{1992}{Silver medal at the International Mathematical Olympiads.}
\end{cvchronology}

\cvsection{Administration of Research}
\begin{cvbullets}
\item Organisation of scientific meetings:
\begin{cvitems}
\item conference ``Formal Power Series and Algebraic Combinatorics 2026'', Seattle (Program committee member).
\item conference ``At the crossroads of physics and mathematics: the joy of integrable combinatorics --- a conference in honour of Philippe Di Francesco's 60th birthday'', 24--26 June 2024, Saclay (France).
\item conference ``Integrability, Combinatorics and Representations '19'', 2--6 September 2019, Giens (France).
\item workshop ``Geometric $R$-matrices'', 18--22 December 2017, MATRIX (Australia).
\item program ``New trends in integrable models'', 15 August--25 November 2016, International Institute of Physics (Natal, Brazil).
\item program ``Statistical Mechanics, Integrability and Combinatorics'', 11 May--3 July 2015, and conference ``Lattice Models: Exact Methods and Combinatorics'', 18--22 May 2015, Galileo Galilei Institute (Florence, Italy).
\item conference ``Integrability and Combinatorics '14'', Giens (France), 23--27 June 2014.
\item conference ``Formal Power Series and Algebraic Combinatorics 2013'', Paris (Program committee member).
\item conference ``Random Matrices and Integrable Systems'', Les Houches, 5--9 March 2012.
\item conference ``GranMa '11'', IHP (Paris), 3--6 October 2011.
\item conference ``Geometry and Integrability in Mathematical Physics '08'', CIRM Luminy (France), 15--19 September 2008.
\item conference ``Geometry and Integrability in Mathematical Physics '06'', Moscow (Russia), 15--19 June 2006.
\item workshop ``Combinatorial methods in Physics and Knot Theory'', Moscow (Russia), 21--25 February 2005.
\end{cvitems}

\item Supervision of six PhD students (T.~Fonseca 2007--2010, A.~Garbali 2012--2015, A.~Gunna 2020--2024, A.~Stoyan 2023--2026, Junru Hu 2026--, Jakob Sloan 2026--), of ten post-docs (L.~Cantini 2006--2008, K.~Kalle 2007--2008, M.~Wheeler and D.~Betea 2012--2014, A.~Ponsaing 2013--2014, P.~Finch 2014--2015, I.~Halacheva 2018--2020, J.~Lamers 2020--2021, J.-E.~Bourgine 2022--2023, L.~Cassia 2023--2025). L.~Cantini later received the CNRS Bronze Medal. Twelve of these are still in academia, among whom five have permanent positions; the remaining four successfully transitioned to the private sector. Supervision of seven Master students and three Bachelor students summer internships.
\item Recipient of four Discovery Projects ``Multivariate polynomials: combinatorics and applications'' (2014--2017), ``Matrix product multi-variable polynomials from quantum algebras'' (2019--2023), ``Elliptic Schubert calculus'' (2021--2026), ``Shuffle algebras and vertex models'' (2024--2026), and a Future Fellowship ``From quantum integrable systems to algebraic geometry and combinatorics'' (2016--2020) from the ARC.
\item Recipient of an ERC consolidator grant for the project ``Loop models, Integrability and Combinatorics'', 2011--2016 (budget: 900k\euro).
\item ``Scientist in charge'' of the Marie Curie Research Training Network ``ENRAGE'' (Random Geometry and Random Matrices: From Quantum Gravity to Econophysics), 2005--2009.
\item Coordinator of the ANR program ``GIMP'' (Geometry and Integrability in Mathematical Physics), 2006--2009.
\item Former editor for {\em Communications in Mathematical Physics}; editor for {\em Algebraic Combinatorics}.
\item Referee for numerous journals (Commun. Math. Phys., Lett. Math. Phys., Electronic J.\ of Combinatorics, Adv.\ Appl.\ Math., J. of Combin. Theory A, JSTAT, J.\ of Physics A, etc); expert for several research organisations (ANR -- France, NWO -- Netherlands, ESF -- Europe, ARC -- Australia).
\item Member of the Committee of Academic Sponsors of SLMath (Berkeley). Former member of the Diversity and Inclusion committee (Melbourne University). Former member of the laboratory council (LPTMS and LPTHE).
\item PhD referee (A.~Ponsaing, L.~Poulain d'Andecy).
\item Member of the European Marie Curie Research Training Network ``ENIGMA'', 2005--2008, of the ESF program ``MISGAM'', 2005--2009, and of the ANR program ``GranMa'', 2009--2011.
\end{cvbullets}

\cvsection{Summary of Publications}
\begin{cvcompact}
\item 80 publications in major international peer-reviewed journals.
\item 7 chapters of books.
\item 8 preprints.
\item 1 popular science article.
\item ``Habilitation \`a Diriger des Recherches'' thesis was published as a monograph.
\item h-index: 33 (Google Scholar, April 2026)
\end{cvcompact}

\cvsection{Summary of Conferences}
\begin{cvcompact}
\begingroup
\newcount\confno
\def\s#1\par{\advance\confno by 1}%
\long\def\newsec#1#2#3\endsec{\confno=0#3\item #2: \the\confno}%
\input conf.tex
\endgroup
\end{cvcompact}

\cvsection{Selected Invited and Plenary Conference Talks}
\begin{cvcompact}
\item Conference ``FPSAC'', UC Davis (USA), 17--21 July 2023: {\sl Schubert puzzles as exactly solvable models} (plenary talk).
\item International Congress of Mathematical Physics, Montreal (Canada), 23--28 July 2018: {\sl From quantum integrability to Schubert calculus}.
\item Conference ``Integrability and Algebraic Combinatorics'', IPAM UCLA (USA), 15--19 April 2024: {\sl The MacNeille completion of the type B Bruhat poset}.
\item Conference ``Quantum integrability and quantum Schubert calculus'', Kavli Royal Society Centre (UK), 11--13 June 2018: {\sl Schubert puzzles and quantum integrability}.
\item Symposium FoPM Tokyo (Japan), 6--8 February 2023: {\sl From exactly solvable models of statistical mechanics to combinatorics}.
\item Main speaker at the 76th ``S\'eminaire Lotharingien de Combinatoire'', 3--6 April 2016: {\sl Razumov--Stroganov Correspondences and the Geometry of Schubert Varieties}.
\end{cvcompact}

\cvsection{Programming, Web Presence}
\begin{cvbullets}
\item Professional expertise in algorithm design.
\item High level programming skills in:
\begin{cvitems}
\item Traditional programming languages: python, c, c++, C\#, Lisp, assembly language\dots;
\item Web-based languages and related: HTML/CSS/JavaScript/TypeScript, PHP/SQL;
\item Mathematics-related languages: Macaulay2, Mathematica.
\end{cvitems}
\item Contributes to various open source projects.\\
Developer for (since 2025, ``author of'') the open source software \href{https://www.macaulay2.com/}{``Macaulay2''}:
\begin{cvitems}
\item Bug fixes, implementation of new functions.
\item Rewrite of the output engine.
\item \href{https://github.com/pzinn/Macaulay2Web}{Macaulay2Web}, a \href{https://www.unimelb-macaulay2.cloud.edu.au}{web-based interface} using the \href{https://katex.org/}{\KaTeX} technology.
\end{cvitems}
\item Contributes entries to the \href{https://oeis.org/}{On-Line Encyclopedia of Integer Sequences}.
\item Various projects can be found on my \href{http://blogs.unimelb.edu.au/paul-zinn-justin/}{website} and on my \href{https://github.com/pzinn/}{GitHub page}.
\end{cvbullets}
